%%%%%%%%%%%%%%%%%%%%%%%%%%%%%%%%%%%%%%%%%%%%%%%%%%%%%%%%%%
%%%%%%%%%%%%%%%%%%%%%%  Question#01  %%%%%%%%%%%%%%%%%%%%%
%%%%%%%%%%%%%%%%%%%%%%%%%%%%%%%%%%%%%%%%%%%%%%%%%%%%%%%%%%
\begin{comment}
  Put all the topics for this question.
\end{comment}
%%%%%%%%%%%%%%%%%%%%%%%%%%%%%%%%%%%%%%%%%%%%%%%%%%%%%%%%%%
\question 
\begin{parts}
%%%%%%%%%%%%%%%%%%%%%%%%%%%%%%%%%%%%%%%%%%%%%%%%%%%%%%%%%%	
\part[2] For units, you can use SI units; \textcolor{blue}{for example}, \textcolor{red}{\qty{2}{\micro\meter}} can be achieved by using the command \verb|\qty{2}{\micro\meter}| or \verb|\SI{2}{\micro\meter}|. This type \verb|\qty{2}{\micro\meter}| is the modern standard in \textcolor{red}{\texttt{siunitx} v3}, while \verb|\SI{2}{\micro\meter}| is the legacy syntax from older versions.



\part[1] In case you need to use the unit for an expression \textcolor{blue}{as for example} \textcolor{red}{ \(20sin(\omega t)\)~\si{\milli\ampere}}, right after the expression, you can use the command \verb|\si{\milli\ampere}|.


\begin{solution}

\end{solution}
%%%%%%%%%%%%%%%%%%%%%%%%%%%%%%%%%%%%%%%%%%%%%%%%%%%%%%%%%%
\part[4] Two resistors of values \qty{1}{\kilo\ohm} and \qty{4}{\kilo\ohm} are connected in series across a constant voltage supply of \qty{100}{\volt}. A voltmeter having an internal resistance of \qty{12}{\kilo\ohm} is connected across the \qty{4}{\kilo\ohm} resistor. Draw the circuit and calculate: 

\begin{subparts}
		\subpart True voltage across \qty{4}{\kilo\ohm} resistor before the voltmeter was connected.
		\subpart Actual voltage across \qty{4}{\kilo\ohm} resistor after the voltmeter is connected and voltage recorded by the voltmeter.
	\end{subparts}



\begin{solution}

\end{solution}
%%%%%%%%%%%%%%%%%%%%%%%%%%%%%%%%%%%%%%%%%%%%%%%%%%%%%%%%%%	
\part[5] \color{red}Figure(pdf, png, etc) can be inserted using the \verb|\includegraphics{nameOfFigure}| command as shown below. : \color{black}Find the rms value of the current waveform of Fig.\ref{fig:currentWaveform}. If the current flows through a \qty{9}{\ohm} resistor, calculate the average power absorbed by the resistor. 


		\begin{figure}[H]
			 \centering 
			 \includegraphics[width=0.30\textwidth]{Fig11.15_ElectricCircuit_SadikuPg446_3E}
			 \caption{}
             \label{fig:currentWaveform}
        \end{figure}
		


\begin{solution}

\end{solution}
%%%%%%%%%%%%%%%%%%%%%%%%%%%%%%%%%%%%%%%%%%%%%%%%%%%%%%%%%%			
\end{parts}

